% page 175 du fac-simile (image p-174 du scan)
% bloc 165.1 x 250.6 mm ; marge gauche 36.9 mm
\pgc{15.240mm}{0.000mm}{
\l{\cel{55}165}
\l{}
\l{}
\l{\sou{fileto}. I. Spiralo qua cirkondas skrubo. - II. Parto di la}
\l{\cel{3}animalo, qua, pro havar nek tendini nek osti, esas distran-}
\l{\cel{3}chebla a lonchi dina. - DFIS.}
\l{}
\l{\sou{filetika}.\cel{2}Qua apartenas a, o relatas, sive la ancestraro di}
\l{\cel{3}speco, sive la procedi qui abutis a la formaceso di la speci,}
\l{\cel{3}sive la diversa lineaji di la grupi di la klasifiko naturala.}
\l{}
\l{\sou{filharmonio}. Amo pasionoza pri muziko. - DEFIRS.}
\l{}
\l{\sou{filheleno}. I. Partisano di la Greki antiqua, di la arti e di la}
\l{\cel{3}civilizeso-sistemo di Grekia antiqua. - II. Amiko di la Greki}
\l{\cel{3}moderna, partisano di lia nedependo politikala. - DEFIS.}
\l{}
\l{\sou{filio}. Persono konsiderata pri la ligilo qua unionas lu a ti de}
\l{\cel{3}qui lu naskis. - EFIRS.}
\l{}
\l{\sou{filialo}. Firmo fondita da firmo \textquotedbl{}patra\textquotedbl{}.}
\l{}
\l{\sou{filigrano}. Laboruro ek fili de oro, de arjento, de vitro, inter-}
\l{\cel{3}plektita ed interligita per solduri preske neperceptebla.}
\l{\cel{3}- DEFIRS.}
\l{}
\l{\sou{filiko}. (\sou{bot.}) Planto senkotiledona, kriptogama, herbacea, di}
\l{\cel{3}qua la semineti kontenesas da kapsuli an la surfaco reversa di}
\l{\cel{3}la folio, di qua la folii volvesas krosatre ante lia deskloz-}
\l{\cel{3}eso, e di qua ula speci divenas arboratra en la regioni tro-}
\l{\cel{3}pikala. - L.\cel{1}\sou{filicea}. - IL.}
\l{}
\l{\sou{filipiko}. Diskurso violentoza kontre ulu qua guvernas la Stato.}
\l{\cel{3}- DEFIRS.}
\l{}
\l{\sou{filistro}. Nomo donata da la kulturozi, da la artisti, a la per-}
\l{\cel{3}soni quin li konsideras kom nekapabla komprenar to quo esas}
\l{\cel{3}spiritala, artala.}
\l{}
\l{\sou{filmo}. Bendo fotografala celuloida, uzata en la kameri obskura}
\l{\cel{3}di la cinematografilo. - DEF.}
\l{}
\l{\sou{filologo}. La cienco qua studias ta o ca linguo nacionala kom}
\l{\cel{3}organo di la vivo mentala, spiritala, di la populo qua uzas}
\l{\cel{3}olu. - DEFIRS.}
\l{}
\l{\sou{filoxero}. (\sou{zool.}) Insekto qua rodas la radiko di la vito. -}
\l{\cel{3}L.\cel{1}\sou{phylloxera vastatrix}. - DEFIRS.}
\l{}
\l{\sou{filozofo.}\cel{1}La persono qua su konsakras a la cienco pri la origini}
\l{\cel{3}e la kauzi.}
\l{}
\l{\sou{filozofiar}. (\sou{netrans.)}\cel{1}I. Rezonar pri la origini e la kauzi. -}
\l{\cel{3}II. Rezonar ciencoze pri ulo. - DEFIRS.}
\l{}
\l{\sou{filtrar.}\cel{1}(\sou{trans.)}\cel{1}Pasigar (liquido) tra korpo poroza (felto,}
\l{\cel{3}stofo, papero, karbono, petro sponjatra, e c.) por purigar lu.}
\l{\cel{3}- DEFIRS.}
}
